% =============================================
% BYLAWS OF THE RENSSELAER UNION STUDENT SENATE
% !TEX root = main.tex
% !LW recipe=latexmk (lualatex)
% Compile with: latexmk main.tex
% =============================================
\documentclass{uniondoc}

\orgname{Rensselaer Union Student Senate}
\orgshort{Student Senate}
\logopath{../Primary_Logo_Red.png}
\amendeddate{April 7, 2026}

\begin{document}

\article{Organization}

\begin{enumerate}
\item The name of this organization shall be the Rensselaer Union Student Senate, hereinafter referred to as the Student Senate or the Senate.
\item The Student Senate gains its authority from the Rensselaer Union Constitution.
\item The Student Senate shall pass legislation commensurate with its duties as defined by the Rensselaer Union Constitution, and as specified in these Bylaws.
\end{enumerate}

\article{Memberships}
\begin{enumerate}
\item A Senator shall be defined as any person who has voting rights in the Senate, as defined by the Rensselaer Union Constitution.
\item The Grand Marshal, Officers of the Senate, and Senators comprise the membership of the Senate.
\item The voting membership of the Student Senate shall consist of all Senator positions that are currently occupied.
\item No voting member of the Student Senate shall simultaneously hold a position on the Rensselaer Union Judicial Board, a voting position on the Rensselaer Union Executive Board, or an elected position on the Rensselaer Union Undergraduate Council.
\item Serving on the Student Senate shall be restricted to Activity Fee-paying members of the Union.
  \begin{enumerate}
  \item For this purpose, students not participating in the Arch summer semester shall be considered as Activity Fee-paying members of the Union from the conclusion of the spring semester until the commencement of the fall semester, provided that they met the requirement of Activity Fee-paying membership during the semesters before and after the summer.
  \end{enumerate}
\end{enumerate}

\article{Officers}

\begin{enumerate}
\item The Grand Marshal is the presiding officer of the Senate.
\item The Officers of the Senate shall be:
  \begin{enumerate}
  \item The chairs of any committee formed by these Bylaws or the Student Senate,
  \item The Vice Grand Marshal for External Affairs,
  \item The Vice Grand Marshal for Internal Affairs,
  \item The Secretary,
  \item The Treasurer,
  \item The Parliamentarian, and
  \item The Senate-Executive Board Liaison.
  \end{enumerate}
\item All Officers of the Senate, with the exception of the Grand Marshal, shall be nominated by the Grand Marshal and confirmed by a 2/3 vote of the Senate.
\item All Officers of the Senate, with the exception of the Grand Marshal and the Senate-Executive Board Liaison, shall serve at the pleasure of the Grand Marshal.
\item The Grand Marshal assumes the duties of any unfilled officer position, with the exception of committee chairs, which the Vice Grand Marshal for External Affairs shall assume the duties of. Neither the Grand Marshal nor the Vice Grand Marshal for External Affairs shall assume the duties of the Elections Commissioner.
  \begin{enumerate}
  \item These duties can be delegated by the Grand Marshal as Interim Officer positions, subject to confirmation by a majority vote of the Senate.
  \item Upon the beginning of a new term of the Senate, all Officers at the conclusion of the past term, except the Grand Marshal, shall serve as Interim Officers by default until an appointment to that position is made and confirmed by the Student Senate, they are relieved of their duties, they step down, or they cease to be members of the Union, whichever comes first.
  \end{enumerate}
\item The Vice Grand Marshal for External Affairs shall be responsible for:
  \begin{enumerate}
  \item Overseeing the Senate’s role in campus advocacy, becoming familiar with the advocacy work of all Senate committees, and serving as an assistant to committee chairs and members on advocacy work;
  \item Meeting regularly with the Grand Marshal and Vice Grand Marshal for Internal Affairs to discuss Senate issues;
  \item Serving, as may be necessary, as a liaison to members of the Institute administration to set up meetings and coordinate joint advocacy projects between committees;
  \item Serving as a substitute for the Grand Marshal at meetings with Institute administrators as may be requested by the Grand Marshal;
  \item Providing guidance and assistance in special operations to raise awareness of issues of importance to students, in collaboration with relevant bodies of the Senate;
  \item Assisting individual Senators that have been charged with investigating issues raised by petition; and
  \item Creating and maintaining documentation to strengthen committee advocacy work.
  \end{enumerate}
\item The Vice Grand Marshal for Internal Affairs shall be responsible for:
  \begin{enumerate}
  \item Serving ex-officio as Chair of the Rules and Administration Committee;
  \item Organizing training and assisting in the onboarding of new officers;
  \item Publicizing meeting times;
  \item Meeting regularly with the Grand Marshal and Vice Grand Marshal for External Affairs to discuss Senate issues;
  \item Acting as the substitute for the Grand Marshal at any meetings of the Senate that the Grand Marshal cannot attend. Should the Vice Grand Marshal for Internal Affairs hold a voting position, at such a meeting, they shall not count towards quorum and may only vote to break a tie;
  \item Overseeing relations between the Senate and other organizations in the Union, including the Union administrative staff; and
  \item Ensuring engagement among Senators and continuity between terms of the Senate.
  \end{enumerate}
\item The Secretary shall be responsible for:
  \begin{enumerate}
  \item Taking meeting attendance and minutes;
  \item Distributing minutes, committee reports, motions, and other printed material to the members of the Student Senate within 48 hours following the adjournment of the meeting; and
  \item Maintaining a roster of current membership of the Senate, along with contact information for each member.
  \end{enumerate}
\item The Treasurer shall be responsible for:
  \begin{enumerate}
  \item Managing funds allocated to the Student Senate as directed by the Grand Marshal or the Student Senate;
  \item Preparing the annual budget of the Student Senate, subject to the approval of the Grand Marshal; and
    \begin{enumerate}
    \item The Grand Marshal or either Vice Grand Marshal must be present at the final budget meeting with the treasurer and the Executive Board Club Representative for the Student Senate.
    \end{enumerate}
  \item Approving transactions conducted by the Senate or any of its Committees. The Grand Marshal may overrule any decision pertaining to such transactions.
  \end{enumerate}
\item The Parliamentarian shall be responsible for:
  \begin{enumerate}
  \item Advising the presiding officer on matters of parliamentary procedure;
  \item Maintaining an orderly meeting, abiding by rules of decorum; and
  \item Maintaining an accurate record of all legislation brought before the Senate.
  \end{enumerate}
\item The Student Senate’s representative on the Rensselaer Union Executive Board, hereinafter referred to as the Executive Board, shall be known as the Senate-Executive Board Liaison. They shall be responsible for:
  \begin{enumerate}
  \item Expressing the Senate’s opinions to the Executive Board and the Executive Board’s opinions to the Senate; and
  \item Reviewing all legislation pertinent to Union fiscal policy and student fees.
  \end{enumerate}
\item The chair of any Senate committee may not necessarily be a voting member of the Senate, but shall regularly attend Senate meetings.
\item The Grand Marshal may, at their discretion, nominate other non-voting officers to assist with administering the business of the Senate.
  \begin{enumerate}
  \item These discretionary non-voting officers shall serve at the pleasure of the Grand Marshal.
  \end{enumerate}
\end{enumerate}

\article{Duties of Office}

\begin{enumerate}
\item Undergraduate Senators, with the exception of the Greek Senators, shall be expected to either chair one Senate committee or sit on at least two Senate committees. Graduate Senators shall be expected to either chair one Senate Committee or sit on at least one Senate committee. Greek Senators shall not be expected to serve on any Senate committee.
\item Senators shall be responsible to dutifully attend all Senate Meetings, unless they obtain an excuse from the Grand Marshal. Senators shall actively contribute to the Student Senate, and shall not hinder its performance.
\item Senators shall be expected to maintain satisfactory attendance in the council that they represent. Satisfactory attendance shall be defined by each council’s bylaws.
\item Failure to meet any of the duties of office shall subject the offending Senator to the removal process as outlined in Article VIII of these Bylaws.
\item It is the responsibility of Senators to follow proper procedures regarding acquiring excuses for absences.
  \begin{enumerate}
  \item Senators shall inform the Grand Marshal of any Senate general meetings they will miss.
  \item Senators shall follow the procedures prescribed in the bylaws of organizations they represent regarding acquiring excuses for absences from those organizations.
  \end{enumerate}
\item It is the responsibility of the Grand Marshal to follow proper procedures regarding Senator absences.
  \begin{enumerate}
  \item The Grand Marshal must inform a Senator if an unexcused absence is noted.
  \item The Grand Marshal may use their discretion in granting absences for upcoming meetings.
  \item The Grand Marshal may consider granting an excuse brought forth after the absence of a Senate meeting if the absence resulted from a sudden circumstance which prohibited them from being able to communicate their absence in a timely manner.
  \end{enumerate}
\end{enumerate}

\article{Meetings}

\begin{enumerate}
\item The Student Senate shall meet at least twice a month during the fall and spring semesters and at least once a month during the summer semester, if there is any new business to attend to. The time and location of these meetings shall be at the discretion of the Grand Marshal.
\item All meetings of the Student Senate shall be open to the public.
\item The Senate may, by a 2/3 vote, close a meeting. Within a closed meeting the following guidelines must be met:
  \begin{enumerate}
  \item All motions, resolutions, and votes passed during the closed meeting must be made public immediately after the meeting is opened; and
  \item The Senate may invite non-Senate members to remain by a majority vote.
  \end{enumerate}
\item The Senate may, by a majority vote, open a closed meeting.
  \item A quorum shall consist of 2/3 of the voting membership of the Student Senate.
        \begin{enumerate}
          \item Business, such as reports and the discussion of motions, may occur without a quorum.
          \item In the absence of a quorum, in addition to the motions allowed in Robert’s Rules of Order, Senators may move and vote on motions to table, postpone, or make any motion related to debate, and restrict who may be present in the meeting.  
          \item During online summer meetings, if a quorum is not present, the Senate may adopt motions using electronic voting as described in Article VII, Clause 3.
        \end{enumerate}
\item The Student Senate shall use the current edition of Roberts Rules of Order as their parliamentary authority for any procedures not specified in these Bylaws.
\item The order of the agenda for a meeting of the Senate may be specified by the meeting’s presiding officer prior to the start of the meeting; if unspecified, the agenda order shall be set by what is specified in Roberts Rules of Order.
\item Meetings may be called by one-quarter of the voting membership of the Senate or by the Grand Marshal. The requesting Senators may name a Senator other than the Grand Marshal to preside over the meeting. The Grand Marshal must be notified of any meeting.
  \begin{enumerate}
  \item The presiding officer is responsible for informing all Senators at least twenty-four hours prior to the meeting as to the date, time, place, and purpose of the meeting. Normal procedures for running this meeting must be followed.
  \end{enumerate}
\end{enumerate}

\article{Legislation}

\begin{enumerate}
\item Legislation shall be defined as any change in the bylaws of the Rensselaer Union, the Union Constitution, the policy of the Rensselaer Union, or official opinion of the Union. Legislation may be introduced for consideration by the Senate in any meeting through the following ways:
  \begin{enumerate}
  \item As a proposal, report, or recommendation of a standing committee or subcommittee,
  \item As a proposal from the Graduate Council or the Undergraduate Council, or
  \item As a proposal sponsored by one or more Senators.
  \item As a proposal from an individual member of the Union, followed by a two-thirds vote of the Senate to consider, after which point it shall be added to the conclusion of the agenda.
  \item As a proposal from an Institute-affiliated and student-led organization, as voted on by simple majority of that organization’s membership, followed by a two-thirds vote of the Senate to consider, after which point it shall be added to the conclusion of the agenda.
  \end{enumerate}
\item Legislation shall be submitted to the Grand Marshal in writing no later than three days prior to the meeting during which it is to be motioned and discussed. Legislation that is proposed by two or more Senators may be submitted and discussed within the same meeting.
\item It shall be the responsibility of the Senate, and primarily the Senators, to initiate all legislation in those areas which affect the Union, including those described by the purpose of the Union.
\item All legislation, unless otherwise specified in the Senate Bylaws or Union Constitution, shall require a majority vote of the Senate to adopt.
  \begin{enumerate}
  \item It shall require a 2/3 vote to determine the amount of the Union Activity Fee. It shall require a 2/3 vote to approve any bylaws submitted to it by another body, as required in the Union Constitution. The Senate may require another body to resubmit its bylaws upon passage of other legislation.
  \item The Student Senate Bylaws may be temporarily suspended by a 2/3 vote. This suspension will stand until any of the following occur: The meeting in which the Bylaws were temporarily suspended adjourned, or the Senate reinstates these Bylaws by majority vote.
  \end{enumerate}
\item Resolutions shall be defined as legislation passing through the Senate solely consisting of a formal expression of opinion on behalf of the Union.
  \begin{enumerate}
  \item The enacting phrase of such legislation shall be ``The Rensselaer Union Student Senate resolves that''
  \item This phrase shall also be used for all procedural matters in the Senate.
  \end{enumerate}
\item Union policy shall be defined as legislation passed through the Senate required to be implemented within the Union,
  \begin{enumerate}
  \item The enacting phrase of such legislation shall be the ``It shall hereby be the policy of the Rensselaer Union''
  \item The Grand Marshal may assign any proposed Union policy to one of the following areas:
    \begin{enumerate}
    \item Pecuniary policy; Mandates on the prioritization of Union funds, which does not include specific budgetary outcomes. This may also include policy for audits and financial transparency of the Union. All legislation in this area shall require a 1/2 vote of the Executive Board.
      \begin{enumerate}
      \item This shall not include the setting of the Union Activity Fee.
      \item Policy of the Senate may require the Executive Board to sever financial relationships if they conflict with other valid legislation of the Union that is non-fiscal in nature. The Senate shall not extend this requirement without a 2/3 vote of its voting membership. Should the Executive Board contest the severing, the Judicial Board shall determine whether the severing necessarily intrudes on the budgetary role of the Executive Board.
      \end{enumerate}
    \item Operational policy; policies setting operational standards for Union facilities, and the broad usage of Union resources, as well as broad goals and priorities. All legislation in this area may be vetoed by the PU, and overridden by a 2/3 vote of the Senate.
    \item Governance Standards; legislation setting standards for the governance of the Union. This may include requiring transparency reports, conflict of interests policy, or reporting requirements. This will also include all rules on selection and appointment, and the Union Media Statement.
      \begin{enumerate}
      \item The Union Media statement will require a 2/3 vote of the Senate to modify.
      \end{enumerate}
    \item Jurisdictional and Limitations Policy; where not provided for in the Union Constitution, the Senate shall adopt legislation establishing the jurisdiction of Union bodies and organizations, and the limits on authority of Union bodies. The Elections Commission will advise if these proposed policies would contradict the Union Constitution, and inform the Senate accordingly.
    \end{enumerate}
  \item The Senate may override any assignment of any such category with a 2/3 vote.
  \item The Senate will not pass any Union policy contradicting any previous Union bylaw.
  \item The Grand Marshal or designated official shall distribute a copy of all legislation passed by the Student Senate to all Officers of the Union and the Director of the Union within 24 hours of its passage and shall also describe the right of any Union body to vote or veto the legislation, as appropriate.
    \begin{enumerate}
    \item Any category of Union policy in which a Union body must vote or can veto legislation shall have two business meetings to take such actions.
      \begin{enumerate}
      \item Should a Union body fail to vote or veto the legislation within the prescribed period, the Union body shall waive their right to do so.
      \end{enumerate}
    \end{enumerate}
  \end{enumerate}
\item The Student Senate shall keep a record of all legislation that comes into consideration. To ensure accurate record keeping and to provide context for posterity, the Senate shall include any supplementary materials that are relevant to the recorded legislation.
\end{enumerate}

\article{Rules for Voting and Debate}

\begin{enumerate}
\item Upon entering debate, a queue will be established and maintained by the Parliamentarian.
  \begin{enumerate}
  \item The queue determines who has the floor to speak. The Parliamentarian is responsible for acknowledging all those who wish to speak and must announce each speaker before their speaking time begins.
  \item Discussion on any given motion shall be limited to 45 minutes. After which, the queue for discussion on the motion will be closed.
    \begin{enumerate}
    \item The discussion time may be extended, or the time limit may be suspended, in which case the time limits shall not apply to the rest of the discussion, by a simple majority vote of the Senate at any time during the discussion.
    \item The limit on discussion shall be enforced and monitored by the parliamentarian.
    \end{enumerate}
  \item Speaking time defaults to five minutes.
  \item Anyone present in the room has the right to speak, and may speak any number of times.
    \begin{enumerate}
    \item Those who have not spoken will be called upon before those who have spoken.
    \item If the queue is on the topic of a presentation, designated presenters may respond to any questions directed at them. Speaking time limits still apply.
    \end{enumerate}
  \item Queues for any motion will proceed until the motion has left the floor. Other queues may be ended by the presiding officer when empty.
  \item A Senator may move to restrict the queue, requiring a second and majority approval.
    \begin{enumerate}
    \item By default, only members of the Senate and Officers may be added to a restricted queue. Additional individuals to be given queue rights may be specified by the Senator moving to restrict the queue.
    \item If an individual without queue rights wishes to speak during a restricted queue, they will not be called upon until no members with queue rights remain on the queue.
    \end{enumerate}
  \item A Senator may move to close the queue, requiring a second and a 2/3 vote.
    \begin{enumerate}
    \item No additional individuals may be added to a closed queue.
    \end{enumerate}
  \item A Senator may move to reopen a closed or restricted queue, requiring a second and a majority vote.
    \begin{enumerate}
    \item When speaking time is limited, speakers can yield their time in the following ways:
      \begin{enumerate}
      \item Yield to the Chair, forfeiting their time, or
      \item Yield to another speaker, transferring their time to another person in the room. That person may decline to speak, forfeiting their time.
      \end{enumerate}
    \end{enumerate}
  \item A yielded speaker may not yield time further.
  \end{enumerate}
\item All voting shall be by a show of hands, unless a Senator, Officer, Committee Chairperson, or the presiding officer requests a roll call vote or unanimous consent:
  \begin{enumerate}
  \item In a show of hands, the presiding officer will request the number in favor, the number opposed, and finally the number of abstentions, in that order. The presiding officer will then announce the vote, and whether the motion has passed. In the event of a tie, the presiding officer may then announce their vote followed by the outcome;
  \item The Parliamentarian will also carry out a count of votes to verify the announced result of a vote;
  \item If a roll call vote is requested, the Secretary will read the names of the Senators alphabetically by surname.
  \item Each Senator may vote yes, no, or pass. If they pass, their name will be read again after the first completion of the roll call. At this time, the Senator must vote or abstain; and
  \item No vote may be determined, nor any result announced, until all Senators present have voted or abstained.
  \item If unanimous consent is requested, the Chair shall ask if there is a member of the Senate that objects to the adoption of the motion being discussed.
  \item A member of the Senate may request unanimous consent during any phase of the debate of a motion,
    \begin{enumerate}
    \item If there are no objections, debate shall end immediately and the motion shall be passed.
    \item Motions passed through unanimous consent shall be recorded as ``agreed by unanimous consent'' on the motion and in the meeting minutes.
    \end{enumerate}
  \end{enumerate}
  \item During the summer semester, should a quorum be absent, the presiding officer is authorized to use electronic voting to conduct business.
        \begin{enumerate}
          \item The Secretary shall post, in a Senate-private channel, a poll with a “Yes”, “No”, and “Abstain” option. Senators shall have 48 hours to vote in the poll. The vote will be considered valid if, within the 48-hour period, it receives enough total votes to reach the 2/3rd quorum requirement. 
        \end{enumerate}
\end{enumerate}

\article{Removals}

\begin{enumerate}
\item Good cause for removal, as provided in the Rensselaer Union Constitution, shall be defined as follows:
  \begin{enumerate}
  \item Two or more unexcused absences from meetings of the Student Senate.
  \item The failure of a member of the Student Senate to conscientiously fulfill the duties of office, as defined in Article IV of these Bylaws.
  \item Misuse of power for personal gain.
  \item Violation of Rensselaer Union or Institute policy such that the violation could constitute a danger to the safety of person or property within the Rensselaer Union or on the premises of the Institute.
  \item General malfeasance.
  \end{enumerate}
\item Senators with more than three unexcused absences from a committee of which they are a member may be removed from the committee. An absence may be excused by the Committee Chair, whose decision may be overruled by the Grand Marshal.
\item A motion to impeach may be brought forth at the request of the organization which the Senator represents, or at the request of two or more Senators. If a motion to impeach is passed by a 2/3 vote of its total voting membership, a removal hearing and removal vote shall be held at the following general Senate meeting.
\end{enumerate}

\article{Committees}

\begin{enumerate}
\item The Chair of each Student Senate standing committee shall be nominated by the Grand Marshal and confirmed by a majority vote.
  \begin{enumerate}
  \item A Vice Chair may be selected by the Chair to assist in management and communication for the committee.
  \item For each committee project or initiative, a Project Lead shall be selected by the committee’s chair;
    \begin{enumerate}
    \item Each project lead may create and lead a student task force, at the counsel of the committee’s chair, to handle initiatives in their area requiring significant effort;
    \item The committee’s chair shall assume the duties of any unfilled project lead position; and
    \item Project leads shall serve at the pleasure of the committee’s chair.
    \end{enumerate}
  \item If the committee does not have a Vice Chair, an acting Chair should be selected by either the committee’s chair, if possible, or the Grand Marshal to fulfill the Chair’s duties in their absence.
  \end{enumerate}
\item Committee membership shall be open to all members of the Rensselaer Union.
  \begin{enumerate}
  \item Any committee member who is not a member of the Senate must be approved by the Committee Chair before receiving membership in the committee.
  \item Members of a Committee serve at the pleasure of the Committee Chair and of the Grand Marshal, and may be removed at any time for good cause.
  \item Good cause for removal includes, but is not limited to, neglecting the duties of membership, failure to adhere to the committee’s rules of order during meetings if applicable, or failure to adhere to these Bylaws when conducting business in the name of the committee.
    \begin{enumerate}
    \item Members are expected to dutifully attend meetings of the committee and complete committee assignments within reasonable deadlines.
    \end{enumerate}
  \end{enumerate}
\item All committees shall be empowered to create subcommittees as necessary to aid them in performing their functions, and subcommittees shall stand at the pleasure of the committee’s chair. The committee’s chair shall be an ex-officio member of any such created or existing subcommittee.
\item All committees shall meet at least twice per month during the fall and spring semesters to work on projects, investigate pending legislation, and to consider proposals for new legislation. Attendance shall be recorded at all committee meetings and submitted to the Grand Marshal. All committee chairs shall report on any open projects within their committees’ purview at each Senate meeting.
\item Each committee chair shall be held accountable for their committee’s performance.
\item During meetings, each Committee shall follow rules of order determined at the discretion of its Chair.
\item The Elections Commission:
  \begin{enumerate}
  \item Consists of a maximum of ten members:
    \begin{enumerate}
    \item The Elections Commissioner, who shall be considered the chair of the committee for the purposes of these Bylaws, one representative from each of the Graduate Council, Undergraduate Council, Executive Board, Judicial Board, Interfraternity Council, Rensselaer Panhellenic Council, one Independent member\footnote[1]{Independent shall be defined as any student not affiliated with a Greek organization.} appointed by the Chair and approved by a majority vote of the committee, one member-at-large appointed by the Chair and approved by a majority vote of the committee, and one member-at-large from the voting membership of the Student Senate appointed by the Grand Marshal.
    \item Each of these groups shall appoint its own representative by the end of the second full week of the fall semester. Should any of the memberships, with the exception of the Commissioner, not be appointed by the required time or be vacant for more than three weeks, the Commissioner shall appoint a member within two weeks of the body’s failure to do so matching the eligibility criteria of the group normally responsible for making the appointment, to be confirmed by majority vote of the Senate. The group responsible for appointing a person to that membership shall forfeit its right to do so until the following semester.
    \item All members must be Activity Fee paying students. All members, with the exception of the chair, shall serve through the end of the spring semester, until replaced, or until resigned, whichever comes first. The chair shall serve until they resign or are replaced.
    \item Quorum for the Elections Commission shall be 2/3 of the total membership.
    \item A member of the Elections Commission who is a candidate for Grand Marshal, President of the Union or joins a political party or is a candidate assistant, as defined by the Student Senate or the Elections Commission, will automatically be removed from the Elections Commission.
    \end{enumerate}
  \item Its duties shall include the following:
    \begin{enumerate}
    \item It shall be responsible for properly publicizing, conducting and supervising all elections for officers and representatives of the governing bodies of the Union. It shall submit a report of the election results to the Student Senate following the election.
    \item It shall, at the discretion of the Senate, establish rules for, supervise, and investigate any student government election or referendum for the Rensselaer Union.
    \item It shall prepare and submit to the Judicial Board and Senate reports on all contested elections, which it investigates.
    \item It shall prepare and allow for the election of a Full Professor with Tenure to the Faculty Senate Promotion and Tenure Committee during the Grand Marshal Week Election.
    \end{enumerate}
  \item The Commissioner shall appoint a Vice Commissioner from the voting membership of the committee. Duties of the Vice Commissioner shall include assisting the Commissioner in the day-to-day administration of the committee. In the event that the Commissioner is absent or unable to run a meeting the Vice Commissioner shall have the authority to chair a meeting. At such a meeting, quorum shall be 2/3 the voting membership, less the Vice Commissioner. The Vice Commissioner shall not have a vote at any meeting that they chair. In the event of the resignation, removal, or extended absence of the Commissioner, the Vice Commissioner shall assume the duties of the Commissioner, until the Commissioner is able to resume their duties, or until a new Commissioner has been appointed.
  \end{enumerate}
\item The Senate Communications and Engagement Committee, hereby referred to as the SCEC, shall be responsible for promoting the initiatives and activities of the Student Senate.
  \begin{enumerate}
  \item Its chairperson may act as the spokesperson for the Student Senate at the discretion of the Grand Marshal, responsible for press releases and other communication with media outlets.
  \item It shall make use of a variety of communication media to publicize the initiatives and activities of the Student Senate, and advocate for issues pertinent to the student body.
    \begin{enumerate}
    \item To this end, it shall coordinate with the Senate’s standing committees to promote current projects and initiatives.
    \item It shall provide the student body with accurate information on the activities, initiatives, policies, and projects of the Senate, and shall serve as the primary public outlet for the Senate’s student advocacy.
    \item It shall work with Institute administration and leverage the advocacy duties of the Student Senate to foster a positive relationship between administration \& students.
    \end{enumerate}
  \item It shall make efforts to ascertain the concerns and opinions of students regarding matters under discussion in the Student Senate.
    \begin{enumerate}
    \item It may solicit student concerns and opinions through public relations initiatives.
    \item It shall inform other Committees of student concerns pertaining to the purview of each Committee.
    \end{enumerate}
  \item It shall foster, among the student body, informed discussion of issues currently facing the Senate.
    \begin{enumerate}
    \item It shall, when necessary, create and distribute informational media in conjunction with its ongoing public relations initiatives.
    \item It shall create and maintain platforms for student inquiry into the activities of the Committees.
    \item It may coordinate publication of electronic media with the Web Technologies Group when necessary.
    \item It may collaborate with the Marketing and Strategies Committee of the Executive Board on projects when necessary.
    \end{enumerate}
  \item It shall promote cohesion between the Senate and other Union bodies.
    \begin{enumerate}
    \item It shall be the responsibility of the chairpersons, or designated representative, of each Student Senate committee to correspond with the SCEC chairperson or designee(s) on all publicity information.
    \item The Director of Public Relations of the Undergraduate Council, a designee of the Graduate Council, and the chairperson of the Judicial Board may correspond with the SCEC and its chairperson if they wish to use the resources and platforms of the committee to promote the information and initiatives of their respective bodies.
    \item The Chair of the SCEC and their equivalents on the Executive Board, Judicial Board, Undergraduate Council, and Graduate Council, or a designee if there is no such equivalent, shall meet at least monthly during the academic year to ensure cohesion between different Union bodies in their communications and messaging.
    \end{enumerate}
  \end{enumerate}
\item The Academic Affairs Committee shall be responsible for initiatives, legislation, and resolutions pertinent to academic curricula, academic advising, admissions, research affairs, post-graduate success, career and professional development, the Registrar’s office, and student-faculty relations.
  \begin{enumerate}
  \item It shall assess the state of Rensselaer academic courses.
    \begin{enumerate}
    \item To this end, it shall coordinate with the Senate Communications and Engagement Committee during outreach programs to ascertain student opinions and concerns regarding academics.
    \item It shall work with relevant organizations within the Institute to identify areas for improvement regarding academics.
    \end{enumerate}
  \item It shall attempt to maintain a diverse membership that represents various academic departments.
    \begin{enumerate}
    \item It shall attempt to include graduate and undergraduate representation on the committee.
    \end{enumerate}
  \item It shall inform the Senate and the student body of any changes in academics at the Institute.
  \item It shall be solely responsible for selecting all student representatives, including undergraduate and graduate representatives, to the Faculty Senate Curriculum Committee.
  \item It shall be responsible for maintaining regular contact between the Student Senate and the Faculty Senate.
  \end{enumerate}
\item The Union Annual Report Committee shall be a standing committee of the Student Senate and Union Executive Board.
  \begin{enumerate}
  \item Its chairperson shall serve at the pleasure of the Grand Marshal and the President of the Union, and its chairperson shall meet with the Grand Marshal and the President of the Union, or their representatives, at least twice a month over the Fall and Spring semesters to maintain cohesion and communication between their respective bodies.
  \item It shall work to produce the Union Annual Report, which will detail the Union Activity Fee for both graduate and undergraduate students; this report will be presented to the Student Senate for approval.
  \item Its membership shall consist of the Director of the Union, as a non-voting member; the Business Administrator, as a non-voting member; the President of the Union as a non-voting ex-officio member; at least one graduate member of the Student Body; and at least one undergraduate member of the Executive Board.
  \end{enumerate}
\item The Facilities and Services Committee shall be responsible for initiatives and legislation pertinent to the facilities and services of the Institute, notably those under the Administration Division.
  \begin{enumerate}
  \item It shall take steps to gather information on the state of services and facilities.
    \begin{enumerate}
    \item To this end, it shall coordinate with the Senate Communications and Engagement Committee during outreach programs to ascertain student opinions and concerns on the state of services and facilities.
    \item It shall, when necessary, address these concerns and coordinate its efforts with proper Institute personnel.
    \item It shall, from time to time, inspect the overall state of facilities and services and report its findings to the Senate.
    \end{enumerate}
  \item It shall take steps to remedy student concerns regarding facilities and services.
    \begin{enumerate}
    \item To this end, it shall be responsible for contacting proper Institute personnel to discuss practical solutions to issues regarding facilities and services.
    \item It shall attempt to implement these solutions, with the support of the Institute and the Senate.
    \end{enumerate}
  \end{enumerate}
\item The Residential Life and Dining Committee shall be responsible for initiatives, legislation, and resolutions pertinent to the quality and future direction of residential life, including dormitory selection processes and off-campus housing and related affairs, dining services on the RPI campus, and the relations between the administration and Residential and Learning Assistants.
  \begin{enumerate}
  \item It shall take steps to gather information on the state of Rensselaer dining services and residential life.
    \begin{enumerate}
    \item It shall develop, maintain, and advertise tools and methods for students to deliver feedback and input into the quality of dormitories, the housing selection process, campus food, dining halls, and other residential and dining services.
    \item It shall, from time to time, inspect the overall state of on-campus living and dining halls and report its findings to the Senate.
    \end{enumerate}
  \item It shall take steps to remedy student concerns regarding campus dining and housing.
    \begin{enumerate}
    \item It shall meet with Hospitality Services,Residential Life, and Off-Campus Commons on a regular basis, at a time and place coordinated between the Committee Chair and relevant administrators.
    \item It shall pursue, research, and propose dining programs that align with student needs, along with any reforms to residential life that are necessary and proper.
    \item It shall collaborate with Hospitality Services, Residential Life, Off-Campus Commons, and other organizations to initiate solutions to student concerns regarding Rensselaer dining services and on-campus living.
    \end{enumerate}
  \item It shall attempt to maintain a diverse membership that represents a wide range of demographics with regards to residential life and campus dining.
  \item It shall attempt to maintain contact with Residential Assistants and Learning Assistants on campus, and to hear their concerns and to recommend policy to the Senate and relevant administrators as may be applicable.
  \end{enumerate}
\item The Student Life Committee shall be responsible for initiatives and legislation pertinent to student rights, quality of life and well-being, Title IX policy and student safety, the Rensselaer Handbook of Student Rights and Responsibilities, and the Division of Student Experience portfolio, including fraternity and sorority life, along with all other campus advocacy areas not already within the jurisdiction of another committee.
  \begin{enumerate}
  \item It shall take steps to gather information on the state of quality of life, well-being, and the general campus student experience.
    \begin{enumerate}
    \item To this end, it shall coordinate with the Senate Communications and Engagement Committee during outreach programs to ascertain student opinions and concerns within these areas.
    \end{enumerate}
  \item It shall address issues pertaining to student rights and responsibilities as they arise.
    \begin{enumerate}
    \item It shall work with the Dean of Students to discuss pending changes to the Handbook of Student Rights and Responsibilities, or other relevant policies that affect student rights.
    \end{enumerate}
  \item It shall take steps to remedy student concerns regarding student life.
    \begin{enumerate}
    \item It shall advise the administration on issues related to student life, and work with them in addressing these issues.
    \item It shall work and collaborate with the administration and other organizations to initiate new programs that improve the student experience.
    \end{enumerate}
  \end{enumerate}
\item The Web Technologies Group shall be responsible for promoting, implementing, and maintaining technological initiatives of the Student Senate.
  \begin{enumerate}
  \item Its chair shall be responsible for assisting the Grand Marshal and Senate Committees in addressing their technological needs.
  \item It shall design and implement technological solutions to student concerns, collaborating with other committees and organizations as necessary.
  \item It shall contribute to Rensselaer Union homepage development efforts and shall be responsible for maintenance and updates of all other Student Government web tools and services.
  \item It shall be responsible for historical cataloging of documents and information for a digital Student Government archive.
  \item It shall be responsible for investigating the status of technology and digital infrastructure operated by the Union, and recommend policies pertinent to those areas to the Senate.
  \end{enumerate}
\item The Community Relations Committee shall be responsible for initiatives pertinent to outreach and service on campus, in the local community, and in the greater Capital Region. This shall include the maintenance and active formation of relationships between the Student Senate and campus organizations, alumni, and local governments in and around the Capital Region.
\item The Rules and Administration Committee shall be responsible for investigating internal affairs of the Union, including its administration and policies, in addition to maintaining the effectiveness of the Senate as a body.
  \begin{enumerate}
  \item It shall investigate questions of governance structures, maintenance of student government, and effective policy-making, and shall serve as a resource to other Senate committees to this end.
  \item It shall investigate and lead initiatives relating to continuous internal maintenance of the Senate, and, in cooperation with the Vice Grand Marshal for External Affairs, other Senate committees, and other relevant bodies, shall investigate and pursue initiatives aimed to improve the Senate’s effectiveness and accessibility.
  \item It shall oversee the Senate’s compliance with its Bylaws.
  \item It shall investigate and recommend proposals for policies pertinent to the affairs of the Union, including pecuniary policy, within the Senate’s legislative power vested to it by the Union Constitution.
  \item This subclause shall not be construed as infringing upon the power of any committee of the Senate to investigate and recommend policy pertaining to its interest area. However, other committees or bodies within the Union may request the assistance of the Rules and Administration Committee in drafting policy or for input in issues of governance.
  \item It shall be responsible for reviewing the language of all constitutional amendments that come before the Senate for approval.
  \item It shall, as may be necessary, serve as an intermediary between the Senate and the Judicial Board on questions of a constitutional nature, and shall investigate parliamentary questions in cooperation with the Senate Parliamentarian.
  \item It shall conduct a yearly review of Union policy, and provide recommendations to the Senate on ways to improve cohesiveness and clarity.
  \item It shall be the chief custodian of historical Senate records.
  \item It shall be charged with the notification of other bodies of the Union in all changes in Union policy passed by the Senate which are relevant to that particular body’s operations or responsibilities.
  \item It shall investigate the status of Union affairs as it pertains to compliance with Union policy, and inform the Senate of all pertinent information on this subject.
  \end{enumerate}
\end{enumerate}

\article{Petitions}

\begin{enumerate}
\item Petitions shall be a means by which members of the Union may bring concerns, initiatives, or ideas before the Student Senate.
  \begin{enumerate}
  \item Petitions may only be created and signed by members of the Union.
    \begin{enumerate}
    \item Should a signatory lose their membership status during the duration of a petition they have signed, their signature on that petition shall remain valid.
    \end{enumerate}
  \item The Rules and Administration Committee shall maintain rules governing the petition process.
    \begin{enumerate}
    \item Such rules shall be approved by a majority vote of the Senate.
    \end{enumerate}
  \end{enumerate}
\item Should any petition, compliant with rules governing the petition process, reach a threshold of 250 signatures within one full calendar year of its posting, it shall be placed on the agenda at a Senate meeting occurring within the next fifteen academic days, excluding exam periods.
  \begin{enumerate}
  \item The Senate may choose to address petitions at any time before this.
  \item Should the Senate’s term end within fifteen academic days of a petition reaching this threshold, that petition may instead be addressed by the next Senate in the semester of their election.
  \end{enumerate}
\item The Senate may choose to take the following courses of action when addressing a petition:
  \begin{enumerate}
  \item Charge one or more Senators to investigate issues raised by the petition.
  \item Charge a committee to investigate the petition as an initiative.
  \item Vote on a resolution on the issue.
  \item Vote to hold a referendum election on the issue.
  \item Refer a petition to another organization as appropriate.
  \end{enumerate}
\end{enumerate}

\article{Appointments}

\begin{enumerate}
\item The Grand Marshal shall make any appointments not specified by these or the Bylaws of other respective boards.
\item Appointments made by the Grand Marshal must be reported to the Senate at its next scheduled meeting, or will be automatically nullified.
\item If either Greek Senator position becomes vacant, the President of either the Interfraternity Council or the Rensselaer Panhellenic Council, for their council’s respective Greek Senator, shall appoint a replacement. The appointment shall be confirmed by simple majority vote of the Student Senate.
\item If a position of Independent Senator becomes vacant, the Grand Marshal shall appoint a replacement. The appointment shall be confirmed by simple majority vote of the Student Senate.
\item The Grand Marshal shall nominate student members of the Review Board as student positions on the Review Board become available, which shall be confirmed by a majority vote of the Senate.
\end{enumerate}

\article{Amendments}

\begin{enumerate}
\item The procedure to amend the Bylaws shall be as follows:
  \begin{enumerate}
  \item All proposed amendments to the Student Senate Bylaws must be proposed in their entirety and passed by a 2/3 vote of the Senate.
  \end{enumerate}
\item These Bylaws shall be amended as necessary so as not to be in conflict with outside laws, the Rensselaer Polytechnic Institute Corporate Charter, the Rensselaer Polytechnic Institute Bylaws, the Rensselaer Polytechnic Institute Student Bill of Rights, the Rensselaer Union Constitution, and decisions of the Board of Trustees of Rensselaer Polytechnic Institute.
\item These Bylaws shall be modified to reflect an approved amendment by:
  \begin{enumerate}
  \item Inserting new words or clauses into these Bylaws at their appropriate places;
  \item Striking out words or clauses from these Bylaws; or
  \item Striking out and inserting words or clauses simultaneously.
  \end{enumerate}
\end{enumerate}

\article{Succession}

\begin{enumerate}
\item If the position of the Grand Marshal is vacant, or if the Grand Marshal is unable to fulfill the duties or requirements of office, then the Vice Grand Marshal for Internal Affairs shall assume the duties of the Grand Marshal as Acting Chair of the Senate.
\item The Acting Chair of the Senate assumes the duties of Grand Marshal until such a time that a meeting of the Senate can be called. The first and only order of business of such a meeting shall be the selection of a new Grand Marshal.
\item All nominations for the position of Grand Marshal must be motioned and seconded by the members of the Senate. The nomination processes shall be overseen by the Chair of the Judicial Board.
\item In the case of multiple nominees, an appointment for Grand Marshal shall be selected by a series of simple majority votes.
  \begin{enumerate}
  \item If there are two nominees, a vote shall be taken, and the nominee with a simple majority shall be named the appointment for Grand Marshal.
  \item If there are three or more nominees, a series of votes shall be taken.
    \begin{enumerate}
    \item In each vote, the nominee with the least number of votes shall be removed from consideration, and only the remaining nominees shall be considered in the next vote.
    \item This series of votes shall continue until a single nominee remains, and that nominee shall be named the appointment for Grand Marshal.
    \end{enumerate}
  \end{enumerate}
\item The nominee who is named the appointment for Grand Marshal must be approved by 2/3 of the Senate’s total voting membership. If the appointment is not approved, a new process of nominations begins.
\item Once the appointment of a new Grand Marshal has been approved by the Senate, the Acting Chair of the Senate shall immediately relinquish the responsibilities of office to the new Grand Marshal.
\end{enumerate}

\article{Inquiries}

\begin{enumerate}
\item The Senate may open an inquiry into areas concerning Union membership and Union affairs.
  \begin{enumerate}
  \item Any legislation opening an inquiry must state the intent of the inquiry and relate it to a particular and legitimate legislative interest of the Senate.
    \begin{enumerate}
    \item In such a case where the subject of an inquiry are the actions of an individual, such an individual must be removable by the Student Senate, any Officer thereof, or be removable by no other body but the Union membership or a constituency thereof as defined in the Union Constitution, and the relevant legislation must specify removal as an interest of inquiry.
    \item No inquiry may be opened into a subject area reserved exclusively to another Union body in the Constitution.
    \end{enumerate}
  \item Inquiries to be delegated to a standing committee, investigatory committee, or individual may be established through a simple majority vote, while inquiries of the whole are to be established through a 2/3 vote.
  \end{enumerate}
\item Inquiries of the whole shall be presided over by the Grand Marshal, and shall be held during meetings of the Senate as a whole.
\item If an inquiry has been delegated to an investigatory committee, and the motion establishing such inquiry does not name the chair of the investigatory committee, the Grand Marshal shall appoint such a chair.
  \begin{enumerate}
  \item Chairs of investigatory committees shall not be considered Officers of the Senate, but shall serve at the pleasure of the Grand Marshal.
    \begin{enumerate}
    \item If the area of interest of the inquiry pertains to actions of the Grand Marshal, the chair may not be relieved of their duties by the Grand Marshal.
    \item If the Senate believes the chair of an investigatory committee has failed to meet the duties of office outlined in Article VIII of these Bylaws, then they may remove the chair by simple majority vote.
    \item If an investigatory committee does not have a chair, the Grand Marshal may appoint one without need for confirmation by the Senate as a whole, unless an appointment by the Grand Marshal is formally objected to by two or more Senators, in which case the appointment shall follow the standard nomination process.
    \end{enumerate}
  \end{enumerate}
\item In the establishment of an inquiry maintained by an investigatory committee or, when applicable, delegated to a standing committee, the Senate may direct any of the following from the relevant committee:
  \begin{enumerate}
  \item The fixing of the time for the committee to report to the Senate,
  \item The fixing of the membership of the committee, and placing restraints on the actions of individual members of the committee,
  \item The fixing of the meeting time or times of the committee, or
  \item Any further constraints on the action and procedure of the committee not explicitly stated.
  \end{enumerate}
\item An investigatory committee, standing committee charged with conducting an inquiry, or individual tasked with conducting an inquiry shall be empowered, in pursuance to the objectives of the inquiry assigned to them by the Senate:
  \begin{enumerate}
  \item To hold public hearings, unless the nature of the case, specific testimony, or specific documents require that the hearing be closed to members of the general public.
  \item To gather testimony, documents, or comments from individuals or bodies relevant to the inquiry.
    \begin{enumerate}
    \item When testimony, documents, or comments are requested, the individuals or bodies in question must be given reasonable notice and proper accommodations to respond.
    \item A lack of testimony from an individual in an inquiry pertaining to their removal shall not be used by the Senate as cause for removal, nor shall any other sanction be imposed as a result of that individual's lack of testimony.
    \end{enumerate}
  \item To compel the provision of documents belonging to the Union, which includes but is not limited to Bylaws, meeting minutes, committee records, and attendance records, in such cases where an individual or body does not comply with a request for such documents pursuant to subclause (b).
    \begin{enumerate}
    \item No documents may be compelled which are reasonably considered to be primarily the private affairs of individuals. All contests regarding this provision shall be appealed to the Judicial Board.
    \end{enumerate}
  \end{enumerate}
\item The Senate may not remove an individual as a result of an inquiry without the interest of the inquiry being declared as such.
  \begin{enumerate}
  \item If an inquiry has removal as a declared interest, the affected individuals shall be notified appropriately prior to the beginning to the process of the inquiry.
    \begin{enumerate}
    \item Such notification shall include adequate notice of the right to choose not to testify.
    \end{enumerate}
  \end{enumerate}
\item For the purposes of this article, the classification of a document as public or private shall be determined by the body providing the document at the time it is provided.
  \begin{enumerate}
  \item Such a body may request that a particular document remain private. In such a case, or in cases where the document was private at the time of request, the Senate and all subsidiary bodies or individuals shall not make such a document public.
  \end{enumerate}
\item If no time limit is specified for an inquiry by the motion establishing an inquiry, the inquiry shall terminate upon a verbal or written report on the outcome of the inquiry to the Senate.
  \begin{enumerate}
  \item Such a report shall be in the public record, unless the Senate votes to make such a report private by a 2/3 vote.
  \item Such a report shall be approved by all individuals conducting an inquiry by a simple majority vote.
  \end{enumerate}
\item The members of an inquiry may close any proceedings of their inquiry with a simple majority vote.
\end{enumerate}

\end{document}

% Local Variables:
% TeX-engine: luatex
% End:
